% Part: incompleteness
% Chapter: theories-computability
% Section: extensions-of-q-not-decidable

\documentclass[../../../include/open-logic-section]{subfiles}

\begin{document}

\olfileid{inc}{tcp}{cqn}

\olsection{$\Th{Q}$-এর সঙ্গতিপূর্ণ প্রসারণগুলি অনির্ণেয়}

\begin{explain}
মনে করো, কোনো তত্ত্বকে \emph{সঙ্গতিপূর্ণ} বলা হয় যদি কোনো সূত্র~$!A$-এর
জন্য সেটি $!A$ ও $\lnot !A$ উভয়ই প্রমাণ না করে। বিরোধ থেকে যেকোনো কিছু
অনুসৃত হয় বলে অসঙ্গতিপূর্ণ তত্ত্ব তুচ্ছ: প্রতিটি বাক্যই সেখানে প্রমাণযোগ্য।
স্পষ্টত, কোনো তত্ত্ব $\omega$-সঙ্গতিপূর্ণ হলে সেটি সঙ্গতিপূর্ণ। কিন্তু
সঙ্গতিপূর্ণ হওয়া দুর্বলতর শর্ত (অর্থাৎ এমন তত্ত্ব আছে যা সঙ্গতিপূর্ণ, কিন্তু
$\omega$-সঙ্গতিপূর্ণ নয়)। তাই \olref[oqn]{thm:oconsis-q}-এর অনুমানকে
সাধারণ সঙ্গতিতে দুর্বল করে আমরা আরও শক্তিশালী উপপাদ্য পাই।
\end{explain}

\begin{lem}
কোনো ``সার্বজনীন গণনসাধ্য সম্বন্ধ'' নেই। অর্থাৎ এমন কোনো দ্বি-স্থানীয়
গণনসাধ্য সম্বন্ধ~$R(x,y)$ নেই যার নিচের ধর্মটি আছে: $S(y)$ যেকোনো
এক-স্থানীয় গণনসাধ্য সম্বন্ধ হলে এমন কোনো~$k$ আছে যাতে প্রত্যেক~$y$-এর
জন্য $S(y)$ সত্য যদি এবং কেবল যদি $R(k,y)$ সত্য।
\end{lem}

\begin{proof}
ধরি $R(x,y)$ একটি সার্বজনীন গণনসাধ্য সম্বন্ধ। $S(y)$-কে
$\lnot R(y,y)$ সম্বন্ধ হিসেবে নিই। $S(y)$ গণনসাধ্য বলে কোনো~$k$-এর জন্য
$S(y)$ ও $R(k,y)$ সমতুল্য। কিন্তু তাহলে $S(k)$ একই সঙ্গে $R(k,k)$ ও
$\lnot R(k,k)$ উভয়ের সমতুল্য, যা বিরোধ।
\end{proof}


\begin{thm}
ধরি $\Th{T}$ এমন যেকোনো সঙ্গতিপূর্ণ তত্ত্ব যা $\Th{Q}$-কে অন্তর্ভুক্ত
করে। তাহলে $\Th{T}$ অনির্ণেয়।
\end{thm}


\begin{proof}
ধরি $\Th{T}$ হলো $\Th{Q}$-এর একটি সঙ্গতিপূর্ণ, নির্ণেয় প্রসারণ।
$\Th{T}$ ব্যবহার করে একটি সার্বজনীন গণনসাধ্য সম্বন্ধ সংজ্ঞায়িত করে আমরা
বিরোধে পৌঁছাব।

$R(x,y)$ ঠিক তখনই সত্য হোক যখন
\begin{quote}
$x$ কোনো সূত্র~$!D(u)$-কে সংকেতায়িত করে এবং $\Th{T}$
$!D(\num y)$ প্রমাণ করে।
\end{quote}
$\Th{T}$-কে নির্ণেয় ধরা হয়েছে বলে $R$ গণনসাধ্য। দেখাই যে $R$
সার্বজনীন। $S(y)$ যেকোনো গণনসাধ্য সম্বন্ধ হলে $\Th{Q}$-তে (এবং তাই
$\Th{T}$-তেও) কোনো সূত্র~$!D_S(u)$ সেটিকে উপস্থাপন করে। তখন প্রত্যেক~$n$-এর
জন্য
\begin{eqnarray*}
S(n) & \lif & T \vdash !D_S(\num n) \\
& \lif & R(\Gn{!D_S(u)}, n)
\end{eqnarray*}
এবং
\begin{eqnarray*}
\lnot S(n) & \lif & T \vdash \lnot !D_S(\num n) \\
& \lif & T \not\vdash !D_S(\num n) \quad \text{($\Th{T}$
  সঙ্গতিপূর্ণ বলে)} \\
& \lif & \lnot R(\Gn{!D_S(u)}, n).
\end{eqnarray*}
অর্থাৎ প্রত্যেক~$y$-এর জন্য $S(y)$ সত্য যদি এবং কেবল যদি
$R(\Gn{!D_S(u)},y)$ সত্য। সুতরাং $R$ সার্বজনীন, এবং আমরা অভীষ্ট বিরোধে
পৌঁছেছি।
% BN-SRC-236 TARGET-CORRECTION-BEGIN
\par\smallskip\noindent\textnormal{\small\textbf{উৎস-সংশোধন (BN-SRC-236):}
হিমায়িত দুই সমীকরণে মেটা-সম্বন্ধ~$S$-এর যুক্তি হিসেবে বস্তু-ভাষার
সংখ্যাপদ~$\num n$ লেখা হয়েছে, যদিও সংজ্ঞায় এবং পরের সারিতে~$R$-এর
যুক্তি স্বাভাবিক সংখ্যা~$n$। বাংলা সমীকরণে উভয় স্থানে $S(n)$ রাখা হয়েছে;
উপস্থাপক সূত্রের ভিতরে $\num n$ অপরিবর্তিত আছে।}
% BN-SRC-236 TARGET-CORRECTION-END
\end{proof}

``সত্য পাটিগণিত'' বলতে তত্ত্ব
$\Setabs{!A}{\Sat{N}{!A}}$-কে বোঝায়; অর্থাৎ পাটিগণিতের ভাষার সেই সব বাক্যের
সেট, যেগুলি মানক ব্যাখ্যায় সত্য।
% BN-SRC-237 TARGET-CORRECTION-BEGIN
\par\smallskip\noindent\textnormal{\small\textbf{উৎস-সংশোধন (BN-SRC-237):}
হিমায়িত সেট-সংজ্ঞায় পরিতৃপ্তি-চিহ্নের মডেল স্থানে স্বাভাবিক সংখ্যার
সেট~$\Nat$ লেখা হয়েছে। এই অংশের ঘোষিত ``মানক ব্যাখ্যা'' এবং আগের
\olref[inc][int][def]{def:standard-model}-এর সঙ্গে মিলিয়ে বাংলা সূত্রে
মানক গঠন~$\Struct{N}$ নির্দেশক $N$ রাখা হয়েছে।}
% BN-SRC-237 TARGET-CORRECTION-END

\begin{cor}
সত্য পাটিগণিত নির্ণেয় নয়।
\end{cor}

%In Section~\olref{undefinability:truth:section} we will state a stronger
%result, due to Tarski.

\end{document}
